\documentclass[10pt]{article}
\usepackage[utf8]{inputenc}
\usepackage[T1]{fontenc}
\usepackage{amsmath}
\usepackage{amsfonts}
\usepackage{amssymb}
\usepackage[version=4]{mhchem}
\usepackage{stmaryrd}
\usepackage{bbold}
\usepackage{graphicx}
\usepackage[export]{adjustbox}
\graphicspath{ {./images/} }
\usepackage{caption}

\begin{document}
\captionsetup{singlelinecheck=false}
Basic graphs

\section*{Edgeworth box}
\begin{itemize}
  \item It is a graphic representation of an exchange economy; no production.
  \item 2 goods: $L=2, \ell=1,2$; commodity space is $\mathbb{R}^{2}$.
  \item 2 consumers: $I=2, i=1,2$.
  \item Consumption sets $X_{i}=\mathbb{R}_{+}^{2}$,
  \item Preferences $\succeq_{i}$ complete, reflexive, transitive, continuous, convex, strongly monotone.
  \item $\omega_{i}=\left(\omega_{i 1}, \omega_{i 2}\right) \in \mathbb{R}_{+}^{2}, \omega_{i \ell}$ is $i$ 's endowment of good $\ell$.
  \item Aggregate endowment $\bar{\omega}=\omega_{1}+\omega_{2}$.
  \item Allocation $\left(x_{1}, x_{2}\right) \in \mathbb{R}_{+}^{4}$ to graph in 2 dimensions.
\end{itemize}

\section*{Graphing an Edgeworth box}
\includegraphics[max width=\textwidth, alt={}, center]{99655b76-8fa5-4fed-86b2-fc92a1963771-04_1563_1284_350_227}\\
\includegraphics[max width=\textwidth, alt={}, center]{99655b76-8fa5-4fed-86b2-fc92a1963771-04_1566_1235_349_1554}

Elements of an Edgeworth box

Edgeworth box

\section*{Allocations}
\begin{itemize}
  \item A (feasible) allocation is an element $x$ in $\left(\mathbb{R}_{+}^{2}\right)^{2}$ with $x_{1}+x_{2}=\bar{\omega}$.
  \item It is sometimes called an exact allocation as opposed to requiring $x_{1} +x_{2} \leq \bar{\omega}$.
  \item Every exact allocation is a point in the Edgeworth Box and vice versa.
  \item A price $p \in \mathbb{R}_{+}^{2}$ is represented as a hyperplane in $\mathbb{R}^{2}$, a straight line.
  \item Budget set: $B_{i}\left(p, \omega_{i}\right)=\left\{x \in \mathbb{R}_{+}^{2}: p \cdot x \leq p \cdot \omega_{i}\right\}$. Note that $B_{i}\left(p, \omega_{i}\right)$ may extend outside the Edgeworth box.
\end{itemize}

\section*{Pareto optimality-Planner's problem}
\begin{center}
\includegraphics[max width=\textwidth, alt={}]{99655b76-8fa5-4fed-86b2-fc92a1963771-08_1350_2463_347_155}
\end{center}

\section*{Graph explanation}
The previous graph illustrates the following points.

\begin{itemize}
  \item Interior optimal allocation without differentiable indifference curves, a kink Pareto optima.
  \item Interior Pareto optimal allocation with differentiable preferences.
  \item Corner Pareto optimal allocation; the allocation is at the boundary of player $j$ 's consumption set. This may occur with or without differentiable preferences.
  \item With differentialbe preferences, first order conditions characterize Pareto optimality at an interior point, see Mas-Colell, Whinston, and Green (1995), 16.F, page 561.
\end{itemize}

\section*{Pareto optimality in the Edgeworth box}
\begin{itemize}
  \item If preferences are smooth, interior Pareto optimal allocations may be computed by equating Marginal Rates of Substitutions $M R S_{1} =M R S_{2}$.
  \item If a PO allocation lies on the boundary of the Edgeworth box, tangency may fail, i.e. $M R S_{1} \neq M R S_{2}$. This is common as agents typically consume a small fraction of available goods.
  \item If preferences are not smooth, Pareto optima allocations need not be points of tangency, i.e. $M R S_{1} \neq M R S_{2}$.
\end{itemize}

\section*{Edgeworth box, decentralization}
\begin{center}
\includegraphics[max width=\textwidth, alt={}]{99655b76-8fa5-4fed-86b2-fc92a1963771-11_1450_1948_350_164}
\end{center}

\section*{Edgeworth box, i's and j's problem}
\begin{center}
\includegraphics[max width=\textwidth, alt={}]{99655b76-8fa5-4fed-86b2-fc92a1963771-12_1450_1953_350_162}
\end{center}

\section*{Edgeworth box, WE}
\begin{center}
\includegraphics[max width=\textwidth, alt={}]{99655b76-8fa5-4fed-86b2-fc92a1963771-13_1462_1955_345_162}
\end{center}

Robinson's economy

\section*{Robinson's economy}
\begin{itemize}
  \item Graphical representation of production economies.
  \item One firm, one agent (Robinson), and two goods, leisure $x_{1}$ and yams $x_{2}$.
  \item Robinson's consumption set:\\
$X_{R}=\left\{\left(x_{1}, x_{2}\right):(\forall i) x_{i} \geq 0, x_{1} \leq 24\right\}$.
  \item Robinson's endowment: $\omega=\left(24, \omega_{2}\right)$.
  \item Preferences are complete, transitive, continuous, strongly monotone, strictly convex, (and as differentiable as convenient).
\end{itemize}

\section*{Robinson's firm, assumptions}
\begin{itemize}
  \item The single firm has technology $Y$.
  \item Robinson is the sole owner of the firm.
  \item $0 \in Y$. Inactivity (i.e. no production) is possible.
  \item $Y+\mathbb{R}_{-}^{L} \subset Y$. Free disposal.
  \item $\mathbb{R}_{+}^{L} \cap Y=\{0\}$. No free production.
  \item $Y$ is convex. DRS or CRS.
  \item $y \in Y, y \neq 0 \Longrightarrow-y \notin Y$. Irreversibility.
  \item And sufficiently differentiable when convenient.
\end{itemize}

Robinson's economy, building the graph\\
\includegraphics[max width=\textwidth, alt={}, center]{99655b76-8fa5-4fed-86b2-fc92a1963771-18_1967_2569_40_154}

\section*{Robinson, planner's problem 1}
\begin{center}
\includegraphics[max width=\textwidth, alt={}]{99655b76-8fa5-4fed-86b2-fc92a1963771-19_1436_2496_357_150}
\end{center}

\section*{Robinson, planner's problem 2}
\begin{center}
\includegraphics[max width=\textwidth, alt={}]{99655b76-8fa5-4fed-86b2-fc92a1963771-20_1452_2576_341_150}
\end{center}

\section*{Robinson, planner's problem 3}
\begin{center}
\includegraphics[max width=\textwidth, alt={}]{99655b76-8fa5-4fed-86b2-fc92a1963771-21_1449_2595_335_135}
\end{center}

\section*{Decentralized problem}
\begin{figure}[h]
\begin{center}
\captionsetup{labelformat=empty}
\caption{The firm}
  \includegraphics[alt={},max width=\textwidth]{99655b76-8fa5-4fed-86b2-fc92a1963771-22_1323_1205_412_157}
\end{center}
\end{figure}

\begin{figure}[h]
\begin{center}
\captionsetup{labelformat=empty}
\caption{The consumer}
  \includegraphics[alt={},max width=\textwidth]{99655b76-8fa5-4fed-86b2-fc92a1963771-22_1323_1338_418_1392}
\end{center}
\end{figure}

\section*{The firm's problem}
$$
\pi(p)=\max _{y \in Y} p \cdot y=\max _{y \in Y} p_{2} y_{2}+p_{1} y_{1}=\max _{y \in Y} p_{2} y_{2}-p_{1}\left|y_{1}\right|
$$

The isoprofit line corresponding to $\pi(p)$ is

$$
\begin{gathered}
\left\{y \in \mathbb{R}^{2}: \pi(p)=p \cdot y\right\} \text { or equivalently } \\
\left\{\left(y_{1}, y_{2}\right) \in \mathbb{R}^{2}: y_{2}=\frac{\pi(p)}{p_{2}}+\frac{p_{1}}{p_{2}}\left|y_{1}\right|\right\} \\
y_{2}^{*}=\frac{\pi(p)}{p_{2}}+\frac{p_{1}}{p_{2}}\left|y_{1}^{*}\right|=\pi(p)+p_{1}\left|y_{1}^{*}\right|=\pi(p)=\pi^{*}
\end{gathered}
$$

normalizing $p_{2}=1$ and choosing $y_{1}=0$ in the profit line.

\section*{Decentralized problem and WE}
\begin{figure}[h]
\begin{center}
\captionsetup{labelformat=empty}
\caption{The firm}
  \includegraphics[alt={},max width=\textwidth]{99655b76-8fa5-4fed-86b2-fc92a1963771-24_1329_1168_406_154}
\end{center}
\end{figure}

\begin{figure}[h]
\begin{center}
\captionsetup{labelformat=empty}
\caption{The consumer}
  \includegraphics[alt={},max width=\textwidth]{99655b76-8fa5-4fed-86b2-fc92a1963771-24_1320_1368_418_1358}
\end{center}
\end{figure}

\section*{Not a WE}
\begin{center}
\includegraphics[max width=\textwidth, alt={}]{99655b76-8fa5-4fed-86b2-fc92a1963771-25_1534_2006_348_154}
\end{center}

\section*{A WE}
\begin{center}
\includegraphics[max width=\textwidth, alt={}]{99655b76-8fa5-4fed-86b2-fc92a1963771-26_1578_2016_329_147}
\end{center}

\section*{Remarks}
\begin{itemize}
  \item The illustration of the welfare theorems also provides a method to identify WE; Negishi approach, (Negishi 1960).
  \item If $Y$ is CRS in Robinson's economy, equilibrium prices are solely determined by technology provided the good is produced in equilibrium. CRS firms must earn zero profit.
\end{itemize}

Preferences do not affect prices or distribution of income.

\section*{References}
Mas-Colell, Andreu, Michael Whinston, and Jerry Green. 1995. Microeconomic Theory. New York, Oxford: Oxford University Press.\\
Negishi, T. 1960. "Welfare Economics and the Existence of Equilibrium for a Competitive Economy." Metroeconomica 12: 92-97.


\end{document}